% Part: lambda-calculus
% Chapter: syntax
% Section: alpha

\documentclass[../../../include/open-logic-section]{subfiles}

\begin{document}

\olfileid{lam}{syn}{alp}
\olsection{تبدیل~$\alpha$}

میان~$\lambd[x][x]$ و~$\lambd[y][y]$ چه رابطه‌ای هست؟ هر دو تابع
همانی را بازنمایی می‌کنند. البته از حیث نحوی ترم‌هایی متفاوت‌اند. تفاوت
آن‌ها تنها در نام متغیر مقید است و یکی حاصل «تغییر نام» متغیر مقید در
دیگری است. این عمل \emph{تبدیل~$\alpha$} نام دارد.

\begin{defn}[تغییر متغیر مقید، $\aconvone$]
  اگر ترم~$M$ شامل رخدادی از~$\lambd[x][N]$ باشد، $y \notin
  \FV{N}$، و~$\Subst{N}{y}{x}$ تعریف شده باشد، آنگاه جایگزین‌کردن این
  رخداد با
  \begin{equation*}
    \lambd[y][\Subst{N}{y}{x}]
  \end{equation*}
  که به~$M'$ می‌انجامد، \emph{تغییر متغیر مقید} نامیده و به‌صورت
  $M \redone[\alpha] M'$ نوشته می‌شود.
\end{defn}

\begin{defn}[سازگاری رابطه]
  رابطهٔ~$R$ روی ترم‌ها \emph{سازگار} نامیده می‌شود اگر شرایط زیر را
  برآورده کند:
  \begin{enumerate}
  \item اگر~$R N N'$، آنگاه~$R \lambd[x][N] \lambd[x][N']$
  \item اگر~$R P P'$، آنگاه~$R (PQ) (P'Q)$
  \item اگر~$R Q Q'$، آنگاه~$R (PQ) (PQ')$
  \end{enumerate}
\end{defn}

پس تعریف را چنین بازنویسی می‌کنیم:
\begin{defn}[تغییر متغیر مقید، $\aconvone$]
  \emph{تغییر متغیر مقید} ($\redone[\alpha]$) کوچک‌ترین رابطهٔ
  سازگار روی ترم‌هاست که شرط زیر را برآورده می‌کند:
  \begin{align*}
    &\lambd[x][N] \redone[\alpha] \lambd[y][\Subst{N}{y}{x}] && \text{اگر
      $x \neq y$ و~$y \notin \FV{N}$} \\
    & &&\text{و~$\Subst{N}{y}{x}$ تعریف شده باشد}
  \end{align*}
\end{defn}

«کوچک‌ترین» در اینجا یعنی رابطه فقط زوج‌هایی را در بر می‌گیرد که
سازگاری و شرط افزوده اقتضا می‌کنند، و نه هیچ زوج دیگری را. بنابراین
می‌توان این رابطه را به‌صورت زیر نیز تعریف کرد:

\begin{defn}[تغییر متغیر مقید، $\aconvone$] \ollabel{defn:aconvone}
  \emph{تغییر متغیر مقید} ($\aconvone$) به‌صورت استقرایی چنین تعریف
  می‌شود:
  \begin{enumerate}
  \item اگر~$N \aconvone N'$، آنگاه~$\lambd[x][N] \aconvone
    \lambd[x][N']$ \ollabel{defn:aconvone1}
  \item اگر~$P \aconvone P'$، آنگاه~$(PQ) \aconvone (P'Q)$ \ollabel{defn:aconvone2}
  \item اگر~$Q \aconvone Q'$، آنگاه~$(PQ) \aconvone (PQ')$ \ollabel{defn:aconvone3}
  \item اگر~$x \neq y$، $y \notin \FV{N}$ و~$\Subst{N}{y}{x}$ تعریف شده باشد، آنگاه
    $\lambd[x][N] \redone[\alpha] \lambd[y][\Subst{N}{y}{x}]$.
    \ollabel{defn:aconvone4}
  \end{enumerate}
\end{defn}

این تعریف‌ها هم‌ارزند، اما اثبات را به‌عنوان تمرین واگذار می‌کنیم. از
این پس تعریف استقرایی را به کار خواهیم برد.

\begin{defn}[تبدیل~$\alpha$، $\aconv$]
  \emph{تبدیل~$\alpha$} ($\aconv$) کوچک‌ترین رابطهٔ بازتابی و تعدی
  روی ترم‌هاست که~$\aconvone$ را در بر می‌گیرد.
\end{defn}

همانند بالا، «کوچک‌ترین» یعنی رابطه فقط زوج‌هایی را در بر می‌گیرد که
تعدی و~$\aconvone$ اقتضا می‌کنند؛ این امر به تعریف هم‌ارز زیر می‌انجامد:

\begin{defn}[تبدیل~$\alpha$، $\aconv$] \ollabel{defn:aconv}
  \emph{تبدیل~$\alpha$} ($\aconv$) به‌صورت استقرایی چنین تعریف می‌شود:
  \begin{enumerate}
  \item اگر~$P \aconv Q$ و~$Q \aconv R$، آنگاه~$P \aconv R$.
    \ollabel{defn:aconv1}
  \item اگر~$P \aconvone Q$، آنگاه~$P \aconv Q$. \ollabel{defn:aconv2}
  \item $P \aconv P$. \ollabel{defn:aconv3}
  \end{enumerate}
\end{defn}

\begin{ex}
  ترم~$\lambd[x][f x]$ تبدیل~$\alpha$ به~$\lambd[y][f y]$ می‌شود، و
  برعکس. به‌طور غیرصوری، هر دو تابعی هستند که آرگومانی می‌پذیرند و حاصل
  اعمال~$f$ بر آن آرگومان را، با ارجاع به متغیر محیطی~$f$، بازمی‌گردانند.

  ترم~$\lambd[x][f x]$ تبدیل~$\alpha$ به~$\lambd[x][g x]$ نمی‌شود.
  به‌طور غیرصوری، این دو به‌ترتیب به متغیرهای محیطی~$f$ و~$g$ ارجاع
  می‌دهند، و همین آن‌ها را توابعی متفاوت می‌سازد: در محیط‌هایی که~$f$
  و~$g$ متفاوت‌اند، رفتارشان نیز متفاوت است.
\end{ex}

\begin{prob}
  آیا زوج‌های زیر از ترم‌ها تبدیل‌پذیر به~$\alpha$ هستند؟
  \begin{enumerate}
  \item $\lambd[x][\lambd[y][x]]$ و~$\lambd[y][\lambd[x][y]]$
  \item $\lambd[x][\lambd[y][x]]$ و~$\lambd[c][\lambd[b][a]]$
  \item $\lambd[x][\lambd[y][x]]$ و~$\lambd[c][\lambd[b][a]]$
  \end{enumerate}
\end{prob}

\begin{lem}\ollabel{lem:fv-one}
  اگر~$P \aconvone Q$، آنگاه~$\FV{P} = \FV{Q}$.
\end{lem}

\begin{proof}
  با استقرا بر !!{derivation} رابطهٔ~$P \aconvone Q$.
  \begin{enumerate}
  \item اگر قاعدهٔ آخر~\olref{defn:aconvone4} باشد، $P$ به‌صورت
    $\lambd[x][N]$ و~$Q$ به‌صورت~$\lambd[y][\Subst{N}{y}{x}]$ است،
    که در آن~$x \neq y$، $y \notin \FV{N}$ و~$\Subst{N}{y}{x}$ تعریف
    شده است. بر حسب اینکه~$x \in \FV{N}$ باشد یا نه، حالت‌ها را جدا
    می‌کنیم:
    \begin{enumerate}
    \item اگر~$x \in \FV{N}$، آنگاه:
      \begin{align*}
        \FV{\lambd[y][\Subst{N}{y}{x}]} &=
        \FV{\Subst{N}{y}{x}} \setminus \{y\} \\
        & = ((\FV{N} \setminus \{x\}) \cup \{y\}) \setminus \{y\}
         && \text{بنا بر \olref[sub]{thm:infv}} \\
        & = \FV{N} \setminus \{x\} \\
        & = \FV{\lambd[x][N]}
      \end{align*}
    \item اگر~$x \notin \FV{N}$، آنگاه:
      \begin{align*}
        \FV{\lambd[y][\Subst{N}{y}{x}]}
        & = \FV{\Subst{N}{y}{x}} \setminus \{y\} \\
        & = \FV{N} \setminus \{x\}
         && \text{بنا بر \olref[sub]{thm:notinfv}} \\
        & = \FV{\lambd[x][N]}.
      \end{align*}
    \end{enumerate}
  \item سه حالت دیگر به‌عنوان تمرین واگذار می‌شوند. 
  \end{enumerate}
\end{proof}

\begin{prob}
  برهان~\olref[lam][syn][alp]{lem:fv-one} را کامل کنید.
\end{prob}

\begin{lem}\ollabel{lem:inv}
  اگر~$P \aconvone Q$، آنگاه~$Q \aconvone P$.
\end{lem}

\begin{proof}
  استقرا بر !!{derivation} رابطهٔ~$P \aconvone Q$.
  \begin{enumerate}
  \item اگر قاعدهٔ آخر~\olref{defn:aconvone4} باشد، $P$ به‌صورت
    $\lambd[x][N]$ و~$Q$ به‌صورت~$\lambd[y][\Subst{N}{y}{x}]$ است،
    که در آن~$x \neq y$، $y \notin \FV{N}$ و~$\Subst{N}{y}{x}$ تعریف
    شده است. نخست، داریم~$y \notin
    \FV{\Subst{N}{y}{x}}$، بنا بر~\olref[sub]{thm:clr}. بنا بر
    \olref[sub]{thm:inv}، ترم
    $\Subst{\Subst{N}{y}{x}}{x}{y}$ نه‌تنها تعریف شده، بلکه برابر
    با~$N$ است. سپس بنا بر~\olref{defn:aconvone4} داریم
    $\lambd[y][\Subst{N}{y}{x}]
    \aconvone \lambd[x][\Subst{\Subst{N}{y}{x}}{x}{y}] =
    \lambd[x][N]$.
  \end{enumerate}
\end{proof}

\begin{prob}
  برهان~\olref[lam][syn][alp]{lem:inv} را کامل کنید.
\end{prob}

\begin{thm}
  تبدیل~$\alpha$ رابطه‌ای هم‌ارزی روی ترم‌هاست؛ یعنی بازتابی، متقارن و
  تعدی است.
\end{thm}

\begin{proof}
  \begin{enumerate}
  \item برای هر ترم~$M$، می‌توان~$M$ را با \emph{صفر} تغییر متغیر مقید
    به~$M$ تبدیل کرد.
  \item اگر~$P$ با رشته‌ای از تغییرهای متغیر مقید به~$Q$ تبدیل~$\alpha$
    شود، از~$Q$ می‌توان این تغییرها را (بنا بر~\olref{lem:inv}) به ترتیب
    معکوس وارونه کرد و~$P$ را به دست آورد.
  \item اگر~$P$ با رشته‌ای از تغییرهای متغیر مقید به~$Q$ تبدیل~$\alpha$
    شود و~$Q$ با رشته‌ای دیگر به~$R$، می‌توان با اعمال نخستین رشته و سپس
    دومین رشته، $P$ را به~$R$ تبدیل کرد.
  \end{enumerate}
\end{proof}

از این پس می‌گوییم~$M$ و~$N$ \emph{هم‌ارز به~$\alpha$} هستند،
$M \aeq N$، اگر و تنها اگر~$M$ به~$N$ تبدیل~$\alpha$ شود (که چنان‌که
همین اکنون نشان دادیم، اگر و تنها اگر~$N$ به~$M$ تبدیل~$\alpha$ شود).

\begin{thm}\ollabel{thm:fv}
  اگر~$M \aeq N$، آنگاه~$\FV{M} = \FV{N}$.
\end{thm}

\begin{proof}
  بی‌درنگ از~\olref{lem:fv-one}.
\end{proof}

\begin{lem}\ollabel{lem:sub:R}
  اگر~$R \aeq R'$ و~$\Subst{M}{R}{y}$ تعریف شده باشد، آنگاه
  $\Subst{M}{R'}{y}$ تعریف شده و نسبت به~$\alpha$ با
  $\Subst{M}{R}{y}$ هم‌ارز است.
\end{lem}

\begin{proof}
  تمرین.
\end{proof}

\begin{prob}
  \olref[lam][syn][alp]{lem:sub:R} را ثابت کنید.
\end{prob}

به یاد آورید که در~\olref[sub]{sec}، جای‌گذاری در برخی حالت‌ها تعریف‌نشده
است؛ بااین‌حال، با به‌کارگیری تبدیل~$\alpha$ روی ترم‌ها می‌توان از راه
تغییر نام متغیرهای مقید، جای‌گذاری را همواره تعریف‌شده ساخت. چنان‌که
قضیهٔ زیر نشان می‌دهد، حاصل هم‌ارزی به~$\alpha$ را حفظ می‌کند:

\begin{thm}\ollabel{thm:sub}
  برای هر~$M$، $R$ و~$y$، ترمی~$M'$ وجود دارد چنان‌که~$M \aeq M'$
  و~$\Subst{M'}{R}{y}$ تعریف شده است. افزون بر این، اگر زوج دیگری
  $M'' \aeq M$ و~$R''$ وجود داشته باشد که~$\Subst{M''}{R''}{y}$ تعریف
  شده و~$R'' \aeq R$، آنگاه
  $\Subst{M'}{R}{y} \aeq \Subst{M''}{R''}{y}$.
\end{thm}

\begin{proof}
  با استقرا بر ساخت~$M$:
  \begin{enumerate}
  \item $M$ متغیر~$z$ است: تمرین.
  \item فرض کنید~$M$ به‌صورت~$\lambd[x][N]$ باشد. متغیری~$z$ برگزینید
    که با~$x$ و~$y$ متفاوت باشد و~$z \notin \FV{N}$ و~$z
    \notin \FV{R}$. بنا بر فرض استقرا، ترمی~$N'$ وجود دارد که
    $N' \aeq N$ و~$\Subst{N'}{z}{x}$ تعریف شده است. پس
    $\lambd[x][N] \aeq \lambd[x][N']$ نیز، بنا بر
    \olref{defn:aconvone}\olref{defn:aconvone1}، برقرار است. اکنون
    $\lambd[x][N'] \aeq \lambd[z][\Subst{N'}{z}{x}]$ بنا بر
    \olref{defn:aconvone}\olref{defn:aconvone4}. این کار ممکن
    است، زیرا~$z \ne x$، $z \notin \FV{N'}$ و~$\Subst{N'}{z}{x}$ تعریف
    شده است. سرانجام، $\Subst{\lambd[z][\Subst{N'}{z}{x}]}{R}{y}$
    تعریف شده است، زیرا~$z \neq y$ و~$z \notin \FV{R}$.

    افزون بر این، اگر~$N''$ و~$R''$ دیگری با همان شرایط وجود داشته باشند،
    \begin{multline*}
      \Subst{(\lambd[z][\Subst{N''}{z}{x}])}{R''}{y} =\\
    \begin{aligned}
      &= \lambd[z][\Subst{\Subst{N''}{z}{x}}{R''}{y}] \\
      &= \lambd[z][\Subst{\Subst{N''}{z}{x}}{R}{y}] && \text{بنا بر
        \olref{lem:sub:R}}\\
      &=\lambd[z][\Subst{\Subst{N'}{z}{x}}{R}{y}]
      && \text{بنا بر فرض استقرا}\\
      &=\Subst{(\lambd[z][\Subst{N'}{z}{x}])}{R}{y}
    \end{aligned}
    \end{multline*}
    \item $M$ به‌صورت~$(PQ)$ است: تمرین.
  \end{enumerate}
\end{proof}

\begin{prob}
  برهان~\olref[lam][syn][alp]{thm:sub} را کامل کنید.
\end{prob}

\begin{cor}\ollabel{cor:sub}
  برای هر~$M$، $R$ و~$y$، زوجی از~$M'$ و~$R'$ وجود دارد چنان‌که
  $M' \aeq M$، $R \aeq R'$ و~$\Subst{M'}{R'}{y}$ تعریف شده است. افزون
  بر این، اگر زوج دیگری~$M'' \aeq M$ و~$R''$ وجود داشته باشد که
  $\Subst{M''}{R''}{y}$ تعریف شده باشد، آنگاه
  $\Subst{M'}{R'}{y} \aeq \Subst{M''}{R''}{y}$.
\end{cor}

\begin{proof}
  بی‌درنگ از~\olref{thm:sub}.
\end{proof}

\end{document}
